\clearpage\thispagestyle{empty}
\pdfbookmark[0]{Quellen zu Gauß an Bode (10. Juli 1826)}{bodejul1826-sources}
\begin{center}\textbf{Quellen zu Gauß an Bode (10. Juli 1826)}\end{center}
\noindent Diese KI-Übertragung ins heutige Deutsch folgt Gauß’ Brief vom 10. Juli 1826 in \href{https://gdz.sub.uni-goettingen.de/id/PPN236007467}{den \textit{Werken}, Band~VI} (Göttingen, 1874), S.~454--455. Johann Elert Bodes \textit{Astronomisches Jahrbuch für 1829} (Berlin, 1826), \href{https://books.google.nl/books?id=LxwPAAAAQAAJ&pg=PA144&redir_esc=y}{S.~144} und \href{https://books.google.nl/books?id=LxwPAAAAQAAJ&pg=PA145&redir_esc=y}{S.~145}, nennt das Einsendedatum und druckt den Brief unter der Überschrift \textit{Meridian-Beobachtungen der Pallas und Ceres im Jahr 1826}. Die Links führen zum Digitalisat in Google Books.
\par\smallskip\noindent Jede Tabelle enthält acht Beobachtungen. Die beiden Doppelpunkte und die drei Striche in der Pallas-Tabelle sowie der Doppelpunkt in der ersten Ceres-Zeile sind aus der Werkausgabe übernommen. Auch im Erstdruck ist von einem Stern 10. Größe die Rede. KI-Übertragung und KI-Satz.
\clearpage
